\documentclass[10pt,a4paper]{article}
\usepackage{amsmath,amssymb,amsthm,graphicx}
\usepackage{hyperref}
\title{Catalytic time-space tradeoffs for randomized logspace by recursive stochastic matrix powering}
\author{Edward Pyne}
\date{doi:10.4230/LIPIcs.CCC.2024.4; CC BY 4.0}
\begin{document}
\maketitle
\noindent\textit{Editorial scope.} The counted result is the general catalytic tradeoff, original Theorems3 and27 together. The logarithmic-workspace endpoint, partial-composition theorem, shifted-rounding primitive and all intermediate claims are uncounted support.\par\medskip
\noindent\textit{Allowed standard prerequisites (editorial).}\begin{itemize}
\item Nisan pairwise-independent hash mixing and hybrid-generator bound; canonicalizer specialization of Putterman-Pyne. Original Lemmas 22,29,31 and bibliography [16],[19]; author local proofs/translations of Lemmas 15,21 and Proposition25 included.
\end{itemize}\medskip
% Faithful transcription of CCC 2024 article 4, pp.4:2,4:3,4:6--4:14,4:19.
% Original result numbers are written explicitly; no mathematical changes.
\paragraph{Definition 1.} Let $\mathsf{CTISP}[T(n),S(n),C(n)]$ be the set of languages recognized by catalytic machines that use $O(S(n))$ workspace and $O(C(n))$ catalytic space on inputs of size $n$, and run in time $\operatorname{poly}(T(n))$ in the worst case.
Note that the worst-case runtime must hold over every catalytic tape, as well as every input.
\paragraph{Definition 4.} A read-once branching program (ROBP) $B$ of width $w$ and length $n$ and alphabet $\{0,1\}^t$ is defined by a function $B:[w]\times\{0,1\}^t\to[w]$.\footnote{The standard definition of ROBPs permits the transition function to differ between the layers. However, as we will always be dealing with programs where $w\geq n$, and we are insensitive to polynomial losses in the width, we can assume all transition functions are the same for clarity.} For $x\in(\{0,1\}^t)^n$, define
\[B[i,x]=B[B[\ldots B[B[i,x_1],x_2]\ldots],x_{n-1}],x_n].\]
It is well known that to derandomize $\mathsf{BPL}$, it suffices to estimate $\Pr_{x\leftarrow U_n}[B[1,x]=1]$ up to error $1/3$ for an ROBP $B$ of length $n$, width $n$ and alphabet $\{0,1\}$.
\paragraph{Definition 6 (Mutual ROBP Hitting).} Given a list of branching programs $(B^1,\ldots,B^n)$ each of length and width $n$, such that $\mathbb{E}[B^j(U_n)]\geq1-1/n^2$ for every $j$, produce a fixed $x\in\{0,1\}^n$ such that
\[\bigwedge_{j\in[n]}B^j(x)=1.\]
\paragraph{Theorem 7.} There is a $\mathsf{CTISP}[n,\log n,\log^2 n]$ algorithm that solves the mutual ROBP hitting problem.
\paragraph{Definition 9 (Catalytic Turing Machine [3]).} A Turing machine $M$ is a catalytic machine using time $T(n)$, workspace $S(n)$, and catalytic space $C(n)$ if it has a work tape, a read-only input tape, a write-only output tape, and a catalytic tape $\mathbf{w}$. We require that for every input $x$ with $|x|=n$ and every $\mathbf{w}$, $M^{\mathbf{w}}(x)$ halts in time at most $T(n)$, using at most $S(n)$ cells on the worktape and $C(n)$ cells on $\mathbf{w}$. Moreover, the final configuration of $\mathbf{w}$ must be equal to its initial configuration, for every $x$ and $\mathbf{w}$.
We now define the notion of a catalytic machine that computes a function. We furthermore define the notion of partially computing a function, where on some tapes $\mathbf{w}$ the machine can output a special failure symbol $\bot$.
\paragraph{Definition 10.} For a function $f:\{0,1\}^*\to\{0,1\}^*$, we say a catalytic machine $M$ (catalytically) computes $f$ if for every $x$ and $\mathbf{w}$, $M^{\mathbf{w}}(x)=f(x)$, and at the end of the computation $\mathbf{w}$ is in its original state. We say that $M$ partially (catalytically) computes $f$ if for every $x$ and $\mathbf{w}$, $M^{\mathbf{w}}(x)\in\{\bot,f(x)\}$, and at the end of the computation (no matter the output) $\mathbf{w}$ is in its original state.
Partial catalytic computation is trivial without further restrictions (as $M$ can always output $\bot$), but we require it as an intermediate step in our analyses. We require a further condition on our machines, that they can revert the catalytic tape at any time without the catalytic tape traversing any new configurations:\footnote{There are existing results related to transforming catalytic algorithms into reversible catalytic algorithms [15]. However, they do not appear to maintain worst-case runtime over the catalytic tape, which is crucial for our results.}
\paragraph{Definition 11.} A catalytic machine $M$ is reversible if for every $x$ and initial configuration $\mathbf{w}$, at any point during the execution of $M^{\mathbf{w}}(x)$, the machine can receive an external REVERT signal. Let $P=P(\mathbf{w})$ denote all prior configurations of the catalytic tape during the execution of $M^{\mathbf{w}}(x)$. After this signal, $M$ must reset $\mathbf{w}$ to the original configuration, and moreover every intermediate configuration of $\mathbf{w}$ during this process must lie in $P$. We require any time bound on $M$ to hold even in the case that $M$ is given the REVERT command at an arbitrary point.
\section*{Catalytic Derandomization: author Section 3}
In this section we prove Theorem 2. We state all our results in terms of catalytic algorithms for the stochastic matrix powering problem, as it is easily compatible with the recursive framework we implement later. Recall a nonnegative matrix is stochastic (resp. substochastic) if all row sums are 1 (resp. at most 1). For a set $S$, let $U_S$ be the uniform distribution over $S$, and let $U_n=U_{\{0,1\}^n}$.
\paragraph{Theorem 14.} There is a $\mathsf{CTISP}[n,\log n,\log^2 n]$ algorithm that, given $n$ and a stochastic matrix $M\in[0,1]^{n\times n}$ where each entry is specified with $O(\log n)$ bits of precision, outputs $\widetilde M\in[0,1]^{n\times n}$ such that $\|\widetilde M-M^n\|_1\leq1/n$.
We recall the existence of efficient algorithms which canonicalize (sub)stochastic matrices, essentially reducing the stochastic matrix powering problem to producing a PRG that fools a branching program.
\paragraph{Lemma 15 ([21,19,7]).} There is a constant $c>0$ and a space $O(\log(nw/\varepsilon))$ algorithm which, given $\varepsilon>0$ and $n,w\in\mathbb N$ where $w\geq n$ and a substochastic matrix $M\in[0,1]^{w\times w}$ with $O(\log w)$ bits of precision, returns a branching program
\[B:[(w/\varepsilon)^c]\times\{0,1\}^m\to[(w/\varepsilon)^c]\]
of width $w'=(w/\varepsilon)^c$ and length $m=n\cdot O(\log(w/\varepsilon))$ where $m$ is a power of two. Moreover, letting $\widetilde M\in[0,1]^{w\times w}$ be the (substochastic) matrix where for $i,j\in[w]$ we define\footnote{Note that this truncates the size from $w'$ back to $w$.} $\widetilde M_{i,j}=\Pr_{x\leftarrow U_n}[B[i,x]=j]$, we have $\|\widetilde M-M^n\|_1\leq\varepsilon$.
As this is not the way these results are stated, we provide a translation in Appendix A. We next define the Nisan PRG, and recall several auxiliary lemmas.
\subsection*{The Nisan PRG}
Given a branching program, we first define the larger alphabet program obtained from duplicating each edge:
\paragraph{Definition 16.} For $t\in\mathbb N$, for $B:[w]\times\{0,1\}^n\to[w]$ of width $w$, let $B^t:[w]\times(\{0,1\}^t)^n\to[w]$ be the branching program of length $n$ and width $w$ over alphabet $\{0,1\}^t$ with transition function $B_t[a,y]=B[a,y_1]$, where $y_1$ is the first bit of $y\in\{0,1\}^t$. Note that $B^t$ can be constructed in space $O(\log(tnw))$ given $B$, and furthermore for every $i,j\in[w]$, $\Pr_{x\leftarrow U_n}[B[i,x]=j]=\Pr_{x\leftarrow U_{(\{0,1\}^t)^n}}[B^t[i,x]=j]$.
We recall a pairwise independent hash family with a very efficient description:
\paragraph{Observation 17.} For every $t\in\mathbb N$, there exists a pairwise independent hash family $H:\{0,1\}^t\to\{0,1\}^t$ such that $|H|=2^{2t}$, and $h\in H$ (which we associate with $h\in\{0,1\}^{2t}$) can be evaluated in space $O(t)$.
Given a (hash) function $h:\{0,1\}^t\to\{0,1\}^t$ and a program $B$, we define an operator that applies a single level of the Nisan construction with hash function $h$.
\paragraph{Definition 18.} Given $B^t:[w]\times(\{0,1\}^t)^n\to[w]$ and $h:\{0,1\}^t\to\{0,1\}^t$, let $B^{t,h}:[w]\times(\{0,1\}^t)^{n/2}\to[w]$ be the width $w$, length $n/2$ program with transition function
\[B_{t,h}[a,x]=B_t[B_t[a,x],h(x)].\]
Using a recursive application of hash functions, we can define the Nisan PRG as follows.
\paragraph{Definition 19.} For $(h_1,\ldots,h_\ell)\in H_t$, define $\operatorname{NIS}_{(h_1,\ldots,h_\ell)}:\{0,1\}^t\to\{0,1\}^{t\cdot n}$ inductively as follows. Let $\operatorname{NIS}_0(x)=x_1$, and for $j\in[\ell]$
\[\operatorname{NIS}_{(h_1,\ldots,h_j)}(x)=(\operatorname{NIS}_{(h_1,\ldots,h_{j-1})}(x)\|\operatorname{NIS}_{(h_1,\ldots,h_{j-1})}(h_j(x))).\]
Note that $B[\cdot,\operatorname{NIS}_{(h_1,\ldots,h_\ell)}(\cdot)]$ and $B_{t,(h_1,\ldots,h_\ell)}[\cdot,\cdot]$ (as defined in Definition 18) are equal as functions.
To analyze the Nisan PRG, we define the notion of a hash function being good for composing two functions, and a PRG being good for a function.
\paragraph{Definition 20.} For every $n,w,t\in\mathbb N$ and $\delta>0$ and $f:[w]\times(\{0,1\}^t)^n\to[w]$ and $G:\{0,1\}^t\to(\{0,1\}^t)^n$, we say that $G$ is $\delta$-good for $f$ if for every $i,j\in[w]$,
\[\left|\Pr_{x\leftarrow U_t}[f[i,G(x)]=j]-\Pr_{x\leftarrow U_{(\{0,1\}^t)^n}}[f[i,x]=j]\right|\leq\delta.\]
Moreover, we say $h\in H$ is $\delta$-good (and $\delta$-bad otherwise) for $f:[w]\times\{0,1\}^t\to[w]$ if $G(x)=(x\|h(x))$ is $\delta$-good for $f[i,(x,y)]=f[f[i,x],y]$.
We recall that a random hash function is good with high probability.
\paragraph{Lemma 21 ([16]).} For every $f$, $\Pr_{h\leftarrow H_t}[h\text{ is }\delta\text{-good for }f]\geq1-w^5(1/\delta)^2/2^t$.
(We provide a proof in Appendix A.) Moreover, a hybrid argument establishes the following.
\paragraph{Lemma 22 ([16]).} For every $B^t:[w]\times(\{0,1\}^t)^n\to[w]$ and $\vec h=(h_1,\ldots,h_\ell)$, suppose for every $i\in[\ell]$, $h_i$ is $\delta$-good for $B_{t,h_1,\ldots,h_{i-1}}$. Then $\operatorname{NIS}_{\vec h}$ is $\delta\cdot nw$-good for $B$.
\subsection*{Catalytic Derandomization}
We now state the main result that powers both of our derandomizations.
\paragraph{Theorem 23.} There is a pair of reversible catalytic algorithms $A,D$ that run in workspace $O(\log(nw/\varepsilon))$, catalytic space $O(\log(n)\cdot\log(nw/\varepsilon))$, and time $\operatorname{poly}(nw/\varepsilon)$ and act as follows. Given $\varepsilon>0$ and a length $n=2^\ell$, width $w$ ROBP $B:[w]\times\{0,1\}^n\to[w]$ where $w\geq n$:
\begin{itemize}
\item The machine $A^{\mathbf w}(B)$ outputs $V\in\{0,1\}^{2\ell}$ and $t=O(\log(nw/\varepsilon))$ and sets the catalytic tape to $\mathbf w'$, such that $(\mathbf w',V)$ contains a (read-only) data structure supporting access to hash functions $\vec h=(h_1,\ldots,h_\ell)$ each on $t$ bits, such that $\operatorname{NIS}_{\vec h}$ is $\varepsilon$-good for $B$.
\item The machine $D^{\mathbf w'}(B,V)$ sets the final catalytic tape configuration to $\mathbf w$.
\end{itemize}
For the extension to search problems, we require a version which takes in a list of ROBP, and constructs a set of hash functions that is simultaneously good for all of them. Our proof is for this more general notion, which immediately implies Theorem 23:
\paragraph{Theorem 24.} There is a pair of reversible catalytic algorithms $A,D$ that run in workspace $O(\log(nw/\varepsilon))$, catalytic space $O(\log(n)\cdot\log(nw/\varepsilon))$, and time $\operatorname{poly}(nw/\varepsilon)$ and act as follows. Given $\varepsilon>0$ and a set of $w$ length $n=2^\ell$, width $w$ ROBPs
\[\mathbf B=(B^1,\ldots,B^w)\]
with $B^j:[w]\times\{0,1\}^n\to[w]$ where $w\geq n$:
\begin{itemize}
\item The machine $A^{\mathbf w}(\mathbf B)$ outputs $V\in\{0,1\}^{2\ell}$ and $t=O(\log(nw/\varepsilon))$ and sets the catalytic tape to $\mathbf w'$, such that $(\mathbf w',V)$ contains a (read-only) data structure supporting access to hash functions $\vec h=(h_1,\ldots,h_\ell)$ each on $t$ bits, such that $\operatorname{NIS}_{\vec h}$ is $\varepsilon$-good for $B^j$ for every $j\in[w]$.
\item The machine $D^{\mathbf w'}(\mathbf B,V)$ sets the final catalytic tape configuration to $\mathbf w$.
\end{itemize}
To make our compression and decompression algorithms work, we require that we can determine if a hash function is good for a branching program at a certain level of the Nisan construction, given pointers to the hash functions at the previous levels:
\paragraph{Proposition 25.} There is a space $O(t+\log(w/\delta))$ algorithm that, given $n,w,t\in\mathbb N$ with $w\geq n$ and $\widetilde h\in\{0,1\}^{2t}$ and $B:[w]\times\{0,1\}^n\to[w]$ and (read only) $\mathbf w$ and pointers $p_1,\ldots,p_r$ such that $\mathbf w[p_i,\ldots,p_i+2t]=h_i$ represents a hash function on $t$ bits, returns if $\widetilde h$ is $\delta$-good for $B_{t,(h_1,\ldots,h_r)}$.
We give a proof in Appendix A, as it essentially follows from the argument of [18]. We can then prove the theorem:
\paragraph{Proof of Theorem 24.} We assume without loss of generality that $n,w$ and $1/\varepsilon$ are powers of two. Set
\[t_0=60\log(w/\varepsilon),\qquad t_1=25\log(w/\varepsilon),\qquad\delta=\varepsilon/nw\geq\varepsilon/w^2,\]
and note that we choose $t_1$ large enough such that a good series of hash functions (for all $B^j$ simultaneously) on $t_1$ bits always exists. The algorithm works as follows. First, virtually divide the catalytic tape as:
\[\mathbf w=(\mathbf w_1\|\mathbf w_2\|\cdots\|\mathbf w_{2\ell})\]
where $|\mathbf w_i|=2t_0$, which we think of as initially holding $h:\{0,1\}^{t_0}\to\{0,1\}^{t_0}$. Note that $|\mathbf w|=\ell\cdot4t_0=O(\log(n)\log(nw/\varepsilon))$ as claimed.
Next, initialize $V\in\{0,1,*\}^{2\ell}$ to indicate if each block is compressed, uncompressed, or unprocessed respectively. The first two cases correspond to the following two formats of the block:
\[\mathbf w_i=\begin{cases}h&V_i=1\\(z\|j\|0^{50\log(w/\varepsilon)})&V_i=0.\end{cases}\]
Informally, the first corresponds to block $i$ originally containing a good hash function for $\mathbf B$, and the second corresponds to block $i$ originally containing a bad hash function for $B^j$, which is thus compressible (in fact, $z$ represents a compressed version of the original data). We define notation for the set of blocks in each configuration:
\paragraph{Definition 26.} For $b\in\{0,1\}$, let $I^b(V)\subseteq[2\ell]$ correspond to the indices such that $V_i=b$, and let $S^b(V)=|I^b(V)|$.
Next, we initialize a counter $i=1$ for the current block. We then iterate over $i=\{1,\ldots,2\ell\}$ until $\max\{S^1(V),S^0(V)\}=\ell$.\footnote{If we exit before $i=2\ell$, set the remaining indices of $V$ to an arbitrary value, which we ignore for clarity of presentation.} For each $i$, the algorithm works as follows. Let $\widetilde h=\mathbf w_i$ be the hash function (on $t_0$ bits) obtained from the current block. We then test if $\widetilde h$ is $\delta$-good for
\[f^j=B^j_{t_0,\vec h_p},\qquad\text{for every }j\in[w].\]
Where
\[\vec h_p=(\mathbf w_{I^1(V)_1},\ldots,\mathbf w_{I^1(V)_{S^1(V)}})\]
corresponds to the hash functions on the preceding good blocks, and $B^j_{t_0,\vec h_p}$ is defined as in Definition 18 applied to $B^j$ and $\vec h_p$. As the index set $I^1(V)$ is easy to generate given $V$, this test can be performed in space $O(\log(nw/\varepsilon))$ without modifying the catalytic tape (and hence also in time $\operatorname{poly}(nw/\varepsilon)$), by Proposition 25.
Given the results of this test, we break into cases depending on if $\widetilde h$ is good:
\begin{itemize}
\item If $\widetilde h$ is $\delta$-good for $f^j$ for every $j\in[w]$, set $V_i=1$.
\item If $\widetilde h$ is $\delta$-bad for some $f^j$, set $V_i=0$ and let $j$ be the first index for which this holds. Next, by enumeration over strings $h\in\{0,1\}^{2t_0}$ (which we can do using the workspace), determine the index of $\widetilde h$ in the set
\[\operatorname{BAD}_{i,j}=\{h\in\{0,1\}^{2t}:h\text{ is }\delta\text{-bad for }B^j_{t_0,\vec h_p}\}\]
where we again perform this test using Proposition 25. Letting the index of $\widetilde h$ in this set be $z$, write
\[\mathbf w_i=(z\|j\|0^{50\log(w/\varepsilon)})\]
(we perform this write operation left to right, and will revert it right to left). We denote the final $50\log(w/\varepsilon)$ bits as free space.
Finally, we claim that we can in fact write these quantities in space $|\mathbf w_i|=2t_0$. The index $j$ requires $\log(w)$ bits. Moreover, we have
\begin{align*}
|\operatorname{BAD}_{i,j}|&=2^{2t_0}\cdot\Pr_{h\leftarrow H}[h\text{ is }\delta\text{-bad for }f]\\
&\leq2^{2t_0}\cdot w^5(1/\delta)^2/2^{t_0}\qquad\text{(Lemma 21)}\\
&\leq2^{2t_0}\cdot(w/\varepsilon)^{7-60}.
\end{align*}
And thus $\log|\operatorname{BAD}_{i,j}|\leq2t_0-51\log(w/\varepsilon)=|\mathbf w_i|-51\log(w/\varepsilon)$. Therefore, we can record all required information as claimed.
\end{itemize}
After processing all blocks, we obtain a catalytic tape $\mathbf w'$ and one of two cases:
\begin{itemize}
\item If $S^1(V)=\ell$, there exist $\vec h=(h_1,\ldots,h_\ell)$ corresponding to the hash functions (on $t_0$ bits) in $I^1(V)$, and these functions are easy to recover from $V$.
\item Else, we must have $S^0(V)=\ell$. For $i\in I^0(V)$, let $F_i$ be the $50\log(w/\varepsilon)$ free bits in $\mathbf w_i$. Note that a description of a hash function $h:\{0,1\}^{t_1}\to\{0,1\}^{t_1}$ is of size $|F_i|$. Iterating over $i\in I^0(V)$ in increasing order, we find (via brute force enumeration) a hash function $\widetilde h$ that is $\delta$-good for
\[f^j=B^j_{t_1,\vec h_p},\quad\text{for every }j\in[w].\]
Where
\[\vec h_p=(\mathbf w_{F_{I^0(V)_1}},\ldots,\mathbf w_{F_{I^0(V)_{i-1}}})\]
corresponds to the ($\delta$-good) hash functions stored on the free space in the preceding indices of $I^0(V)$. Next, store $\widetilde h$ in $\mathbf w_{F_i}$. Such a good hash function always exists, by our choice of $t_1$ and Lemma 21, and moreover testing if each candidate is good can be computed in the desired space and time by Proposition 25. After this processing,
\[(\mathbf w_{F_{I^0(V)_1}},\ldots,\mathbf w_{F_{I^0(V)_\ell}})\]
contain $\ell$ good hash functions, which we can clearly access in read-only fashion given given $V$ and $\mathbf w'$.
\end{itemize}
Thus, in both cases we obtain a set of hash functions $\vec h=(h_1,\ldots,h_\ell)$ on $t=O(\log(nw/\varepsilon))$ bits that is $\delta$-good for every one of the relevant tests, so by Lemma 22 we have that $\operatorname{NIS}_{\vec h}$ is $\delta\cdot nw\leq\varepsilon$-good for $B^j$ for every $j$.
\paragraph{Decompression and Reversibility.} It suffices to show that at any point, the algorithm can revert the tape to the original configuration $\mathbf w$ (and then $D(\mathbf B,V)$ simply issues the REVERT command). No matter the present configuration, we iterate through $I^0(V)$ in descending order. Letting the current index be $i\in I^0(V)$, recall this block is of form $(\mathbf w')_i=(z\|j\|*)$. First write $0^{50\log(w/\varepsilon)}$ to the last indices (in reverse order to satisfy reversibility), such that we reach the configuration after compressing the block. Then enumerate over $h\in\{0,1\}^{2t_0}$ using workspace $O(t_0+\log(w/\varepsilon))$, until we find the hash with index $z$ in $\operatorname{BAD}_{i,j}$, where $\operatorname{BAD}_{i,j}$ and $B^j_{t_0,\vec h_p}$ are defined as before (which we still have access to because we reset the tape in reverse order), and we determine membership by Proposition 25.
Once we find this $h$, write $(\mathbf w')_i=h$ (in the reverse order to satisfy reversibility) and proceed to the next highest index in $I^0(V)$. After this process has completed, it is clear from construction that $\mathbf w$ has been reset to the original configuration, and that the tape never reaches a new intermediate configuration during this process.
\paragraph{Time and Space.} In every step of the computation, we perform at most $\operatorname{poly}(2^{t_0}nw/\varepsilon)$ work to determine if a hash function is good, find the index of a bad hash function, or find a good hash function. Moreover, as at every point we store at most a constant number of hash functions on the worktape, the space consumption follows.\hfill$\triangleleft$
It is easy to go from Theorem 23 to Theorem 14.
\paragraph{Proof of Theorem 14.} Let $B:[\operatorname{poly}(n)]\times\{0,1\}^{n^c}\to[\operatorname{poly}(n)]$ be the ROBP obtained from applying Lemma 15 to $M$ with $n=n$, $w=n$, and $\varepsilon=1/2n$. We then call Theorem 23 with $B=B$ and $\varepsilon=1/2n^2$. Let $\vec h=(h_1,\ldots,h_{c\log n})$ be the hash functions obtained from this call, which we have implicit access to via the current state of the catalytic tape $\mathbf w'$ and $V$, and let $t=O(\log n)$ be the domain of the hash functions. Then enumerate over $x\in\{0,1\}^t$ and for $i,j\in[w]$ let
\[\widetilde M_{i,j}=\Pr_{x\leftarrow U_t}[B[i,\operatorname{NIS}_{\vec h}(x)]=j].\]
By Lemma 15 and Theorem 23, we have the guarantee that
\[\|\widetilde M-M^n\|_1\leq\frac1{2n}+n\cdot\frac1{2n^2}=1/n.\]\hfill$\triangleleft$
Finally, we can use Theorem 24 to prove Theorem 7.
\paragraph{Proof of Theorem 7.} The algorithm works as follows. Given $\mathbf B=(B^1,\ldots,B^n)$ of width and length $n$ where $\mathbb E[B^j(U_n)]\geq1-1/n^2$ for every $j$, we call Theorem 24 with $\mathbf B=\mathbf B$ and $\varepsilon=1/2n^2$. Let $\vec h=(h_1,\ldots,h_{\log n})$ be the hash functions obtained from this call, which we have implicit access to via the current state of the catalytic tape $\mathbf w'$ and $V$, and let $t=O(\log n)$ be the length of the domain of the hash functions. We then claim there is some $x\in\{0,1\}^t$ such that for all $j\in[n]$,
\[B^j(\operatorname{NIS}_{\vec h}(x))=1.\]
We have this as
\begin{align*}
\Pr_{x\leftarrow U_t}\left[\bigwedge_{j\in[n]}B^j(\operatorname{NIS}_{\vec h}(x))\right]&\geq1-\sum_{j\in[n]}\Pr_{x\leftarrow U_t}[B^j(\operatorname{NIS}_{\vec h}(x))=0]\\
&\geq1-\sum_{j\in[n]}\left(\Pr_{x\leftarrow U_t}[B^j(U_n)=0]+\frac1{2n^2}\right)\\
&\geq1-n(2/n^2)>0
\end{align*}
where the second inequality follows from our choice of $\varepsilon$. Then it is simple to find such an $x$ by enumeration, whereupon we compute and return $\operatorname{NIS}_{\vec h}(x)\in\{0,1\}^n$, which is precisely the solution for the mutual ROBP hitting problem. Finally, we use Theorem 24 to reset the tape.\hfill$\triangleleft$
\section*{Author Appendix A: proofs used above}
We first prove that a hash function drawn from a pairwise independent hash family is good for a function with high probability. To do this, we recall the hash mixing lemma:
\paragraph{Lemma 29 ([16]).} Let $A,A'\subseteq\{0,1\}^t$ be arbitrary subsets of density $\rho=|A|/2^t$ and $\rho'=|A'|/2^t$. Then for every $\delta>0$,
\[\Pr_{h\leftarrow H}\left[\left|\Pr_{x\leftarrow U_t}[x\in A,h(x)\in A']-\rho\rho'\right|\geq\delta\right]\leq(1/\delta^2)/2^t.\]
\paragraph{Proof of Lemma 21.} For every $i,k\in[w]$, let $A_{i,k}=\{x\in\{0,1\}^t:f[i,x]=k\}$ and let $\rho_{i,k}=|A_{i,k}|/2^t$. Note that for every $i,j$,
\[\Pr_{x,x'\leftarrow U_t}[f[f[i,x],x']=j]=\sum_{k\in[w]}\rho_{i,k}\rho_{k,j}.\]
Thus, for every $h$ such that for every $i,k,j$, $|\Pr_{x\leftarrow U_t}[x\in A_{i,k},h(x)\in A_{k,j}]-\rho_{i,k}\rho_{k,j}|\leq\delta/w$, we have that $h$ is $\delta$-good for $f$. By Lemma 29 this event occurs with probability $1-(w/\delta)^2/2^t$ over $h\leftarrow H$ for each pair of sets, and thus with probability $1-w^3(w/\delta)^2/2^t=1-w^5(1/\delta)^2/2^t$ for every tuple $(i,j,k)$.\hfill$\triangleleft$
We recall there is a logspace algorithm which tests if a hash function is good, given oracle access to the function we wish to fool. We remark that there is work [17,6,20] on testing if an entire PRG is good for a branching program, but we need a much weaker claim.
\paragraph{Lemma 30 ([18]).} There is a space $O(t+\log(w/\delta))$ algorithm that, given oracle access to $f:[w]\times\{0,1\}^t\to[w]$ and $h\in H_t$ and $\delta>0$, tests if $h$ is $\delta$-good for $f$.
\paragraph{Proof.} The algorithm enumerates over $i,j\in[w]$. For every $i,j$, the algorithm computes $p_{i,j}=\mathbb E_{x,x'\leftarrow U_t}[f[f[i,x],x']]$ (i.e. the correct probability) by enumeration over $x,x'$ in space $O(t+\log w)$. Then it computes $\widetilde p_{i,j}=\mathbb E_{x\leftarrow U_t}[f[f[i,x],h(x)]]$ and rejects if the estimate is greater than $\delta$ from the true value. Correctness and total space consumption are immediate.\hfill$\triangleleft$
We can then prove that we can test if a hash function is good, given $B$ and pointers to preceding hash functions.
\paragraph{Proposition 25.} There is a space $O(t+\log(w/\delta))$ algorithm that, given $n,w,t\in\mathbb N$ with $w\geq n$ and $\widetilde h\in\{0,1\}^{2t}$ and $B:[w]\times\{0,1\}^n\to[w]$ and (read only) $\mathbf w$ and pointers $p_1,\ldots,p_r$ such that $\mathbf w[p_i,\ldots,p_i+2t]=h_i$ represents a hash function on $t$ bits, returns if $\widetilde h$ is $\delta$-good for $B_{t,(h_1,\ldots,h_r)}$.
\paragraph{Proof.} By Lemma 30, it suffices to show that given $i\in[w]$ and $x\in\{0,1\}^t$, we can compute $B_{t,(h_1,\ldots,h_r)}[i,x]$. To do this, the algorithm maintains $v\in[w]$ as its current position in the branching program (initialized to $v=i$) and $i\in[n]$ to track the current layer. To determine the next position, it suffices to determine the $i$th block of the output of $\operatorname{NIS}_{(h_1,\ldots,h_r)}(x)$. It is well known that this can be computed in space $O(t+\log nw)$ given read-only access to the set of hash functions (by walking down the binary expansion of $i$, denoted $\langle i\rangle$, and applying $h_j$ if $\langle i\rangle_j=1$), which we have via the pointers.\hfill$\triangleleft$

\paragraph{Author Appendix A: translation of Lemma 15.} Finally, we provide a translation of our quantization statement. We first recall a strict specialization of the statement of [19]:
\paragraph{Lemma 31.} There exists a canonicalizer algorithm $C_t$ that, given $n,w\in\mathbb N$ with $w\geq n$, takes in $\varepsilon>0$ and a sub-stochastic matrix $M\in\mathbb R^{w\times w}$ with each entry represented by at most $O(\log(w/\varepsilon))$ bits, runs in space $O(\log(w/\varepsilon))$, and returns a branching program $B$ of length $n$ and width $w+1$ with alphabet $\{0,1\}^t$ for $t=O(\log(w/\varepsilon))$. Moreover, letting $\widetilde M\in[0,1]^{w\times w}$ be the matrix where for $i,j\in[w]$ we have
\[\widetilde M_{i,j}=\Pr_{x\leftarrow U_{(\{0,1\}^t)^n}}[B[i,x]=j]\]
then
\[\|\widetilde M-M^n\|_1\leq\varepsilon.\]
We reduce the alphabet (and slightly increase the length) as follows. We transform $B$ into a branching program of width $(w+1)\cdot2^t=\operatorname{poly}(w/\varepsilon)$ and length $n\cdot t=n\cdot O(\log(w/\varepsilon))$, where each set of $t$ layers reads $t$ bits, interprets these bits as $\sigma\in\{0,1\}^t$, and takes the transition labeled with $\sigma$ in $B$.

% Faithful transcription of official CCC2024 article4, pp.4:2,4:8--4:9,4:15--4:17.

\paragraph{Theorem 3.} For every $\alpha\in[1,1.5]$,
\[\mathsf{BPL}\subseteq\mathsf{CTISP}[2^{\log^\alpha(n)},\log^\alpha(n),\log^{3-\alpha}(n)].\]

\section*{Author Section 2.1: Composition of Catalytic Algorithms}
We state the main result of this section, which is that catalytic algorithms can be composed without increasing the catalytic space usage. We must be careful when dealing with partial catalytic machines, and in this case we only obtain composition if the machines are reversible (Definition 11).
\paragraph{Theorem 12 (Composition of Partial Catalytic Machines).} Suppose reversible catalytic machines $M_1,M_2$ partially compute $f_1,f_2:\{0,1\}^n\to\{0,1\}^n$ respectively using workspace $S(n)\geq\log(n)$, catalytic space $C(n)$, and time $T(n)$. Then there is a reversible catalytic machine $M$ that partially computes $f_2\circ f_1$ using workspace $2S(n)+O(\log(Sn))$, catalytic space $C$, and time $\operatorname{poly}(T(n))$. Moreover, $M^{\mathbf w}(x)=\bot$ only if $M_2^{\mathbf w}(f_1(x))=\bot$, or there exists $\mathbf w'$ such that $M_1^{\mathbf w'}(x)=\bot$.
\paragraph{Proof.} We proceed roughly following the standard proof for composition of space-bounded algorithms. We maintain two sections on the worktape of size $S$ for $M_1$ and $M_2$ (and auxiliary state of size $O(\log(Sn))$ to keep track of the location of read and write heads, and the FSM configuration of both machines), and a single catalytic tape $\mathbf w$.
\paragraph{The Simulation.} We now begin to simulate $M_2^{\mathbf w}$. We first verify that $M_1^{\mathbf w}(x)\neq\bot$. As in the conventional composition of space-bounded algorithms, every time $M_2$ reads its input, we run $M_1^{\mathbf w}(x)$ on a separate section of the worktape and return the relevant bit of its output, where $\mathbf w$ is the same catalytic tape used by $M_2$, in whatever its current configuration is at the time of the tape read. Moreover, every time $M_2$ writes to the catalytic tape, resulting in a configuration $\mathbf w'$, we run $M_1^{\mathbf w'}(x)$ and verify that it does not produce $\bot$.
\paragraph{Computing the Function.} In the case that $M_1^{\mathbf w'}(x)=f_1(x)$ for every configuration $\mathbf w'$ that is encountered in this simulation and $M_2^{\mathbf w}(f_1(x))=f_2(f_1(x))$, it is easy to see that $M^{\mathbf w}(x)=f_2(f_1(x))$. Moreover, it is clear that in this case we successfully reset the tape.
Otherwise, consider the first point at which $M_1^{\mathbf w'}(x)=\bot$. We first undo the most recent change to the catalytic tape, and send the REVERT command to $M_2$. Once $M_2$ has finished reverting, return $\bot$. We claim that $M_2$ successfully reverts the tape. This follows from the reversibility of $M_2$, and the fact that all calls $M_2$ makes to its input during this process are correctly answered by $M_1$. The latter property follows as every time $M_2$ queries its input during the revert process, $\mathbf w$ is in a state that was encountered during the forward pass, and hence $M_1(x)$ produced $f_1(x)$ when initialized with this catalytic configuration (as otherwise we would have aborted sooner).
\paragraph{Reversibility.} Essentially the same argument establishes that $M$ is reversible. If we receive the REVERT command, let $P(\mathbf w)$ be the states of the catalytic tape that have been encountered so far. First send REVERT to $M_1$ (if operating), and once it has completed send the REVERT command to $M_2$. Moreover, while $M_2$ is reverting, we claim that $\mathbf w$ remains in $P(\mathbf w)$. This follows from the fact that $M_2$ is reversible (as any configurations it creates will lie in $P(\mathbf w)$), and moreover every time $M_2$ queries its input, any computation done by $M_1$ will likewise keep $\mathbf w$ in $P(\mathbf w)$, as we already called $M_1$ with this starting configuration in the forward pass (and $M_1$ will not produce $\bot$, as otherwise we would have already aborted). Thus, the composed algorithm is reversible.
\paragraph{Time and Space.} We now argue the space and time are as claimed. There are a constant number of pointers (which we maintain on the worktape) to track the number of bits output by $M_1$, current tape heads, and other information. The fact that the catalytic tape size is preserved is immediate. The call overhead adds at most a polynomial factor in the runtime, as we run $M_1$ at most once per step of $M_2$. Finally, if $M_1$ computes $f_1(x)$ for every catalytic tape and $M_2$ computes correctly on $\mathbf w$, we successfully compute $f_2\circ f_1$ as claimed.\hfill$\triangleleft$
We derive an easy corollary in the case of multiple composition:
\paragraph{Corollary 13.} Suppose a reversible catalytic machine $M$ partially computes $f:\{0,1\}^n\to\{0,1\}^n$ using workspace $S(n)\geq\log n$, catalytic space $C$, and time $2^S$. Then there exists a reversible catalytic machine $M'$ that partially computes $f^\ell$ using workspace $O(\ell\cdot S)$, catalytic space $C$, and time $2^{O(\ell\cdot S)}$. Moreover, for $x$ where $M^{\mathbf w}(f^i(x))=f(f^i(x))$ for every $\mathbf w$ and $i\in\{0,\ldots,\ell-1\}$, $M'^{\mathbf w}(x)=f^\ell(x)$.
\section*{Author Section 4: Catalytic Recursive Matrix Powering}
We now transform Theorem 23 into a parameterized algorithm for matrix powering.

\paragraph{Theorem 27.} There is a catalytic machine that, given $r_1,r_2$ such that $r_1r_2=\log(n)$ and a stochastic matrix $M\in[0,1]^{n\times n}$ where each entry is specified with $l=O(\log n)$ bits of precision, uses workspace $O(r_2\cdot\log(n))$, catalytic space $O(r_1\cdot\log(n))$, and time $n^{O(r_2)}$, and outputs $\widetilde M$ such that $\|\widetilde M-M^n\|\leq1/n$.

Theorem 27 immediately implies Theorem 3 by setting $r_2=\log^{\alpha-1}(n)$ and $r_1=\log^{2-\alpha}(n)$ for $\alpha\in[1,1.5]$, and using the standard transformation of estimating the acceptance probability of a BPL machine via stochastic matrix powering.
We first prove there exists an algorithm which computes a single intermediate power. We must be careful to ensure that the algorithm satisfies the requirements of (partial) catalytic computation. In particular, if the machine ever outputs an answer (rather than $\bot$), this must be the only possible answer for this input, over all possible catalytic tapes. Simultaneously, we must ensure that the vast majority of inputs never return $\bot$ no matter the initial catalytic tape configuration.
We achieve this dual guarantee using an idea from Saks and Zhou [21]. For a given input matrix $M$, we additionally take in a shift $s\in[0,\delta]$ for $\delta=1/\operatorname{poly}(w/\varepsilon)$. After computing an approximate power of $M$, we add $s$ to each entry, and then truncate each entry to $O(\log(w/\varepsilon))$ bits of precision. In fact, we first verify that our shifted approximate power is sufficiently far from the rounding threshold, and if not return $\bot$. By doing so, we algorithmically verify that we will never round to different thresholds over different $\mathbf w$. Unfortunately, for some pairs $(M,s)$ it may be the case that we detect possible mis-rounding for some tapes $\mathbf w$, even if all possible approximations lie inside a single rounding interval. This can result in returning $\bot$ on some tapes and a (consistent) value otherwise. However, we show that we can choose the magnitude of $s$ such that with high probability over $s$ this does not occur, and we always return a (consistent) value.
\paragraph{Theorem 28.} For every $n,w\in\mathbb N$ and $\varepsilon>0$ where $w\geq n\geq\log w$ and $2^{-n}>\varepsilon>0$, there is a reversible catalytic machine $P$ that uses workspace $O(\log(w/\varepsilon))$, catalytic space $O(\log(n)\log(nw))$, and time $\operatorname{poly}(nw/\varepsilon)$. The machine takes input $s\in\{0,1\}^{O(\log(w/\varepsilon))}$ and a substochastic matrix $M\in[0,1]^{w\times w}$, where each entry of $M$ is specified with $l=O(\log(w/\varepsilon))$ bits of precision. Moreover:
\begin{itemize}
\item For every $(s,M)$, there is substochastic $\widetilde M_s$ (defined without reference to $\mathbf w$) with $l$ bits of precision satisfying $\|\widetilde M_s-M^n\|_1\leq\varepsilon$ such that for every $\mathbf w$, $P^{\mathbf w}(s,M)\in\{\bot,\widetilde M_s\}$.
\item For every $M$, $\Pr_s[\exists\mathbf w,P^{\mathbf w}(s,M)=\bot]\leq1/w^2$.
\end{itemize}
\paragraph{Proof.} Let $B:[(w/\delta)^c]\times\{0,1\}^{m=n\cdot O(\log(w/\delta))}\to[(w/\delta)^c]$ be the result of Lemma 15 applied to $M$ with $n=n$ and $\varepsilon=\delta$ to be chosen later. We compose the output of this algorithm with Theorem 23, applied with $B=B$ and $w'=(w/\delta)^c$ and $n'=m=O(n^3)$ and $\varepsilon=\delta$ to be chosen later. Let $\vec h=(h_1,\ldots,h_\ell)$ be the hash functions obtained from this call, which we have implicit access to via the current state of the catalytic tape $\mathbf w'$ and $V$, and let $t=O(\log(w/\delta))$ be the domain of the hash functions. Then enumerate over $x\in\{0,1\}^t$ and for every $i,j\in[w]$ let
\[\widetilde M_{i,j}=\Pr_{x\leftarrow U_t}[B[i,\operatorname{NIS}_{\vec h}(x)]=j].\]
Next, define $\tau=s\cdot2^{-2k}\cdot J$ where $J=1^{w\times w}$, interpreting $s\in[2^k]$. We next check if any entry of $(\widetilde M+\tau)$ is within $2\delta$ of a multiple of $2^{-l}$, our rounding threshold. In this case, run $D^{\mathbf w'}(B,V)$ to reset the tape and return $\bot$. Otherwise, let
\[\widetilde M_s=\lfloor\widetilde M+\tau\rfloor_l\]
where $\lfloor\cdot\rfloor_l$ rounds each entry down to $l$ bits of precision, and decreases the largest entry per row such that the final matrix is substochastic. Let this matrix be $\widetilde M_s$. Finally, run $D^{\mathbf w'}(B,V)$, and return $\widetilde M_s$.
\paragraph{Accuracy.} By our choice of error in Theorem 23 and Lemma 15, we have that
\[\|\widetilde M-M^n\|_1\leq2w\delta\]
and moreover $\widetilde M$ has each row sum at most 1. Furthermore, perturbing by $\tau$ and rounding down the largest entry causes an $\ell_1$ error of at most $2w\cdot2^{-k}$. Finally, rounding each entry down to a multiple of $2^{-l}$ causes a total error of at most $w\cdot2^{-l}$, so
\[\|\widetilde M_s-M^n\|\leq2w\delta+2w\cdot2^{-k}+w\cdot2^{-l}\leq\varepsilon\]
Where the final inequality comes from choosing
\[l=O(\log(w/\varepsilon)),\qquad k=10\cdot l,\qquad\delta=2^{-2k}.\]
\paragraph{Uniqueness.} We claim that for every $(s,M)$, there is at most 1 possible non-$\bot$ output over all choices of $\mathbf w$ (which we denote $\widetilde M_s$ in the theorem statement). Let $\widetilde M^{\mathbf w},\widetilde M^{\mathbf w'}$ be the result of Theorem 23 on $M$ initialized with catalytic tapes $\mathbf w,\mathbf w'$. By the accuracy guarantee of Theorem 23, for every $i,j$ we have
\begin{align*}
|(\widetilde M^{\mathbf w}+\tau)_{i,j}-(\widetilde M^{\mathbf w'}+\tau)_{i,j}|&\leq|(\widetilde M^{\mathbf w}+\tau)_{i,j}-(M+\tau)_{i,j}|\\
&\quad+|(M+\tau)_{i,j}-(\widetilde M^{\mathbf w'}+\tau)_{i,j}|\leq2\delta.
\end{align*}
Thus, if $(\widetilde M^{\mathbf w'}+\tau)_{i,j}$ is greater than $2\delta$ from a multiple of $2^{-l}$, we can be certain that no tape $\mathbf w'$ will induce an estimate that falls on the other size of the threshold, and hence all non-$\bot$ outputs will be rounded consistently.
\paragraph{Success Probability.} Furthermore, we argue that for every $M$, with probability at least $1-1/w^2$ over $s$ we return $\widetilde M_s$ (not $\bot$) for every initial tape configuration. Fixing arbitrary $s$ and $i,j\in[w]$, if $(M+\tau)_{i,j}$ is at least $3\delta$ from every multiple of $2^{-l}$, every $\mathbf w$ will induce an estimate $(\widetilde M+\tau)_{i,j}$ that is at least $2\delta$ from every multiple of $2^{-l}$, and hence for every $\mathbf w$ we will not produce $\bot$ due to this entry. This occurs for every $i,j$ simultaneously with probability at least
\[w^2\cdot2^l\cdot\frac{6\delta\cdot2^k+2}{2^k}\ll1/w^2.\]
\paragraph{Time and Space.} It is clear the algorithm runs in the claimed time and space bound, given Theorem 23.
\paragraph{Reversibility.} As the only components of the algorithm that write to the catalytic tape are calls to Theorem 23, reversibility follows immediately from the equivalent result for that algorithm.\hfill$\triangleleft$
We can then prove the main result.
\paragraph{Theorem 27 (restated).} There is a catalytic machine that, given $r_1,r_2$ such that $r_1r_2=\log(n)$ and a stochastic matrix $M\in[0,1]^{n\times n}$ where each entry is specified with $l=O(\log n)$ bits of precision, uses workspace $O(r_2\cdot\log(n))$, catalytic space $O(r_1\cdot\log(n))$, and time $n^{O(r_2)}$, and outputs $\widetilde M$ such that $\|\widetilde M-M^n\|\leq1/n$.
\paragraph{Proof.} Let $\vec s=(s_1,\ldots,s_{r_2})\in\{0,1\}^{r_2\cdot O(\log n)}$ be a vector of random shifts. Let $\widetilde M_0=M$ and for $i\in[r_2]$ recursively define
\[\widetilde M_i=f(\widetilde M_{i-1},s_i),\]
where $f$ is the function defined by Theorem 28 with $\varepsilon=1/n^3$ and $n=2^{r_1}$. An easy inductive proof [19] establishes that, letting
\[\|\widetilde M_i-M^{2^{r_1\cdot i}}\|_1=\delta_i\]
we have $\delta_{i+1}\leq1/n^3+2^{r_1}\cdot\delta_{i-1}\leq2^{r_1+1}\cdot\delta_{i-1}$, and hence $\delta_{r_2}\leq1/n$.
The final algorithm iterates over $\vec s$ and computes $\widetilde M_{r_2}$ by applying recursive composition of space-bounded machines Corollary 13 to the algorithm of Theorem 28 as defined above.\footnote{We can define the machine to take the entire shift vector and a pointer to the index it should use, such that we are recursively composing exactly the same machine, but we suppress this for clarity.} The algorithm returns the first non-$\bot$ output. The fact that the algorithm is catalytic follows from Corollary 13 and Theorem 28. Next, we claim there is some $\vec s$ where the algorithm returns a value. Note that $s_i$ is chosen obliviously to $\widetilde M_{i-1}$, and so with probability at least $1/n^2$ over $s_i$, on input $(\widetilde M_{i-1},s_i)$ the algorithm returns $\widetilde M_i$ (i.e. not $\bot$) when run with every possible catalytic tape. Thus, there is some $\vec s$ where every level computes correctly.
\paragraph{Time and Space.} Every application of Theorem 28 occurs with parameters $n=2^{r_1}$ and $w=n$ and $\varepsilon=1/n^3$, such that the algorithm uses workspace $O(r_1+\log(n))=O(\log n)$, catalytic space $O(r_1\cdot(r_1+\log(n)))=O(r_1\cdot\log n)$, and time $\operatorname{poly}(n)$, and moreover the shift $s$ for each level is of length $O(\log n)$. Applying Corollary 13, we obtain that the composed algorithm uses workspace $O(r_2\cdot\log(n)+|\vec s|)=O(r_2\cdot\log n)$, catalytic space $O(r_1\cdot\log n)$, and runs in time $n^{O(r_2)}$ as claimed.\hfill$\triangleleft$

\begin{thebibliography}{22}
\bibitem[1]{ref1} Dimitris Achlioptas. Database-friendly random projections: Johnson-lindenstrauss with binary coins. J. Comput. Syst. Sci., 66(4):671–687, 2003. doi:10.1016/S0022-0000(03)00025-4.
\bibitem[2]{ref2} Sagar Bisoyi, Krishnamoorthy Dinesh, and Jayalal Sarma. On pure space vs catalytic space. Theor. Comput. Sci., 921:112–126, 2022. doi:10.1016/j.tcs.2022.04.005.
\bibitem[3]{ref3} Harry Buhrman, Richard Cleve, Michal Koucký, Bruno Loff, and Florian Speelman. Computing with a full memory: catalytic space. In Symposium on Theory of Computing, STOC 2014, pages 857–866. ACM, 2014.
\bibitem[4]{ref4} Harry Buhrman, Michal Koucký, Bruno Loff, and Florian Speelman. Catalytic space: Nondeterminism and hierarchy. Theory Comput. Syst., 62(1):116–135, 2018. doi:10.1007/ s00224-017-9784-7.
\bibitem[5]{ref5} Jin-yi Cai, Venkatesan T. Chakaravarthy, and Dieter van Melkebeek. Time-space tradeoff in derandomizing probabilistic logspace. Theory Comput. Syst., 39(1):189–208, 2006. doi: 10.1007/S00224-005-1264-9.
\bibitem[6]{ref6} Kuan Cheng and William M. Hoza. Hitting sets give two-sided derandomization of small space. Theory of Computing, 18(21):1–32, 2022. doi:10.4086/toc.2022.v018a021.
\bibitem[7]{ref7} Gil Cohen, Dean Doron, Ori Sberlo, and Amnon Ta-Shma. Approximating iterated multiplication of stochastic matrices in small space. In Proceedings of the 55th Annual ACM Symposium on Theory of Computing, STOC 2023, pages 35–45. ACM, 2023.
\bibitem[8]{ref8} James Cook and Ian Mertz. Catalytic approaches to the tree evaluation problem. In Proccedings of the 52nd Annual ACM SIGACT Symposium on Theory of Computing, STOC 2020, pages 752–760. ACM, 2020.
\bibitem[9]{ref9} James Cook and Ian Mertz. Tree evaluation is in space o(log n · log log n). Electron. Colloquium Comput. Complex., TR23-174, 2023. arXiv:TR23-174.
\bibitem[10]{ref10} Samir Datta, Chetan Gupta, Rahul Jain, Vimal Raj Sharma, and Raghunath Tewari. Randomized and symmetric catalytic computation. In Computer Science – Theory and Applications – 15th International Computer Science Symposium in Russia, CSR 2020, pages 211–223. Springer, 2020.
\bibitem[11]{ref11} Dean Doron, Edward Pyne, and Roei Tell. Opening up the distinguisher: A hardness to randomness approach for BPL = L that uses properties of BPL. Electron. Colloquium Comput. Complex., TR23-208, 2023. arXiv:TR23-208.
\bibitem[12]{ref12} Chetan Gupta, Rahul Jain, Vimal Raj Sharma, and Raghunath Tewari. Unambiguous catalytic computation. In 39th IARCS Annual Conference on Foundations of Software Technology and Theoretical Computer Science, FSTTCS 2019, volume 150 of LIPIcs, pages 16:1–16:13, 2019.
\bibitem[13]{ref13} Daniel M. Kane, Raghu Meka, and Jelani Nelson. Almost optimal explicit johnson-lindenstrauss families. In Leslie Ann Goldberg, Klaus Jansen, R. Ravi, and José D. P. Rolim, editors, Approximation, Randomization, and Combinatorial Optimization. Algorithms and Techniques – 14th International Workshop, APPROX 2011, and 15th International Workshop, RANDOM 2011, Princeton, NJ, USA, August 17-19, 2011. Proceedings, volume 6845 of Lecture Notes in Computer Science, pages 628–639. Springer, 2011. doi:10.1007/978-3-642-22935-0\_53.
\bibitem[14]{ref14} Daniel M. Kane and Jelani Nelson. Sparser johnson-lindenstrauss transforms. J. ACM, 61(1):4:1–4:23, 2014. doi:10.1145/2559902.
\bibitem[15]{ref15} Ian Mertz. Reusing space: Techniques and open problems. Bulletin of EATCS, 141(3), 2023.
\bibitem[16]{ref16} Noam Nisan. Psuedorandom generators for space-bounded computation. In Harriet Ortiz, editor, Proceedings of the 22nd Annual ACM Symposium on Theory of Computing, May 13-17, 1990, Baltimore, Maryland, USA, pages 204–212. ACM, 1990. doi:10.1145/100216.100242.
\bibitem[17]{ref17} Noam Nisan. On read-once vs. multiple access to randomness in logspace. Theor. Comput. Sci., 107(1):135–144, 1993. doi:10.1016/0304-3975(93)90258-U.
\bibitem[18]{ref18} Noam Nisan. RL $\leq$ SC. Comput. Complex., 4:1–11, 1994. doi:10.1007/BF01205052.
\bibitem[19]{ref19} Aaron (Louie) Putterman and Edward Pyne. Near-optimal derandomization of medium-width branching programs. In Proceedings of the 55th Annual ACM Symposium on Theory of Computing, STOC 2023, pages 23–34, 2023.
\bibitem[20]{ref20} Edward Pyne, Ran Raz, and Wei Zhan. Certified hardness vs. randomness for log-space. In 64th IEEE Annual Symposium on Foundations of Computer Science, FOCS 2023, Santa Cruz, CA, USA, November 6-9, 2023, pages 989–1007. IEEE, 2023. doi:10.1109/FOCS57990.2023.00061.
\bibitem[21]{ref21} Michael E. Saks and Shiyu Zhou. $\mathsf{BP}_H\mathsf{Space}(S)\subseteq\mathsf{DSPACE}(S^{3/2})$. J. Comput. Syst. Sci., 58(2):376–403, 1999. doi:10.1006/jcss.1998.1616.
\bibitem[22]{ref22} D. Sivakumar. Algorithmic derandomization via complexity theory. In John H. Reif, editor, Proceedings on 34th Annual ACM Symposium on Theory of Computing, May 19-21, 2002, Montréal, Québec, Canada, pages 619–626. ACM, 2002. doi:10.1145/509907.509996.
\end{thebibliography}
\end{document}
